\documentclass[a4paper,12pt]{article}

% 页面设置
\usepackage[top=2.54cm, bottom=2.54cm, left=1.91cm, right=1.91cm]{geometry}

% 中文支持
\usepackage[UTF8]{ctex}
\usepackage{fontspec}

% 数学公式支持
\usepackage{amsmath}
\usepackage{amssymb}
\usepackage{amsfonts}
\usepackage{amsthm}

\usepackage{enumitem}

% 定理环境定义
\newtheorem*{pf}{证}
\newtheorem*{sol}{解}

% 图形支持
\usepackage{graphicx}
\usepackage{tikz}

% 超链接支持
\usepackage{hyperref}
\hypersetup{
    colorlinks=true,
    linkcolor=blue,
    filecolor=magenta,
    urlcolor=cyan,
}

% 行距设置
\linespread{1.25}

% 标题信息
\title{Mathematical Analysis II\\Week 4 Homework (March 26)}
\author{袁子强\hspace{1cm} 2025533009}
% \date{}

\begin{document}

\maketitle

\section*{Section 9.2}

19. 求下列复合函数的偏导数或导数.

(2) 设 $u = \mathrm{e}^{xyz}, x = rs, y = \frac{r}{s}, z = r^s$, 求 $\frac{\partial u}{\partial r}, \frac{\partial u}{\partial s}$;
\begin{sol}
    将 $x = rs, y = \frac{r}{s}, z = r^s$ 代入 $u = \mathrm{e}^{xyz}$，得
    $$
        u = \mathrm{e}^{r^{s+2}}.
    $$
    因此
    $$
        \frac{\partial u}{\partial r} = (s+2)r^{s+1}\mathrm{e}^{r^{s+2}},
    $$
    $$
        \frac{\partial u}{\partial s} = r^{s+2}\mathrm{e}^{r^{s+2}} \ln r.
    $$
\end{sol}

(4) 设 $u = \frac{\mathrm{e}^{ax}(y-z)}{a^2+1}, y = a\sin x, z = \cos x$, 求 $\frac{\mathrm{d}u}{\mathrm{d}x}$;
\begin{sol}
    将 $y = a\sin x, z = \cos x$ 代入 $u = \frac{\mathrm{e}^{ax}(y-z)}{a^2+1}$，得
    $$
        u = \frac{\mathrm{e}^{ax}(a\sin x-\cos x)}{a^2+1}.
    $$
    因此
    $$
        \frac{\mathrm{d}u}{\mathrm{d}x}
        = \frac{a\mathrm{e}^{ax}(a\sin x-\cos x) + \mathrm{e}^{ax}(a\cos x+\sin x)}{a^2+1}
        = \mathrm{e}^{ax}\sin x.
    $$
\end{sol}


20. 求下列复合函数的偏导数或导数，其中各题中的 $f$ 均有连续的二阶偏导.

(2) 设 $u = f(x,y,z), x = \sin t, y = \cos t, z = \mathrm{e}^t$, 求 $\frac{\mathrm{d}u}{\mathrm{d}t}$;
\begin{sol}
    由链式法则得
    $$
        \frac{\mathrm{d}u}{\mathrm{d}t}
        = \frac{\partial f}{\partial x}\frac{\mathrm{d}x}{\mathrm{d}t}
        + \frac{\partial f}{\partial y}\frac{\mathrm{d}y}{\mathrm{d}t}
        + \frac{\partial f}{\partial z}\frac{\mathrm{d}z}{\mathrm{d}t}
        = f_x \cos t - f_y \sin t + f_z \mathrm{e}^t.
    $$
\end{sol}

(4) 设 $u = f(x + y + z, x^2 + y^2 + z^2)$, 求 $\frac{\partial u}{\partial x}, \frac{\partial^2 u}{\partial x^2}, \frac{\partial^2 u}{\partial x\partial y}$;
\begin{sol}
    令 $f = f(s,t)$，其中 $s = x + y + z, t = x^2 + y^2 + z^2$.

    首先计算一阶偏导数：
    $$
        \frac{\partial u}{\partial x}
        = \frac{\partial f}{\partial s}\frac{\partial s}{\partial x}
        + \frac{\partial f}{\partial t}\frac{\partial t}{\partial x}
        = f_s + 2x f_t.
    $$

    再计算二阶偏导数：
    $$
        \frac{\partial^2 u}{\partial x^2}
        = \frac{\partial}{\partial x}\left(\frac{\partial u}{\partial x}\right)
        = \frac{\partial f_s}{\partial x} + 2f_t + 2x\frac{\partial f_t}{\partial x}
        = f_{ss} + 2x f_{st} + 2f_t + 2x f_{ts} + 4x^2 f_{tt}
        = f_{ss} + 4x f_{st} + 2f_t + 4x^2 f_{tt}.
    $$
    $$
        \frac{\partial^2 u}{\partial x\partial y}
        = \frac{\partial f_s}{\partial y} + 2x\frac{\partial f_t}{\partial y}
        = f_{ss} + 2y f_{st} + 2x f_{ts} + 4xy f_{tt}
        = f_{ss} + 2(x+y)f_{st} + 4xy f_{tt}.
    $$
\end{sol}


22. 试求函数 $z = \arctan \frac{y}{x}$ 在圆 $x^2 + y^2 - 2x = 0$ 上一点 $P\left( \frac{1}{2}, \frac{\sqrt{3}}{2} \right)$ 处沿该圆周逆时针方向上的方向微商.
\begin{sol}
    圆的标准方程为
    $$
        (x-1)^2 + y^2 = 1.
    $$

    因此当 $t = \frac{2\pi}{3}$ 时，对应点 $P\left(\frac{1}{2}, \frac{\sqrt{3}}{2}\right)$.

    首先计算偏导数：
    $$
        \frac{\partial z}{\partial x} = -\frac{y}{x^2+y^2}, \quad
        \frac{\partial z}{\partial y} = \frac{x}{x^2+y^2}.
    $$

    在点 $P$ 处，$x^2+y^2 = 1$，因此梯度为
    $$
        \nabla z|_P = \left(-\frac{\sqrt{3}}{2}, \frac{1}{2}\right).
    $$

    圆的法向量为 $\overrightarrow{n} = (x-1, y)$，因此逆时针切向向量为
    $$
        \overrightarrow{\tau} = (-y, x-1) = \left(-\frac{\sqrt{3}}{2}, -\frac{1}{2}\right).
    $$

    因此方向微商为
    $$
        \frac{\partial z}{\partial \tau}
        = \nabla z \cdot \overrightarrow{\tau}
        = \frac{3}{4} - \frac{1}{4}
        = \frac{1}{2}.
    $$
\end{sol}


23. 求函数 $u = x^2 + 2y^2 + 3z^2 + xy + 3x - 2y - 6z$ 在点 $(1,1,-1)$ 的梯度和最大方向微商.
\begin{sol}
    首先计算一阶偏导数：
    $$
        \frac{\partial u}{\partial x} = 2x + y + 3, \quad
        \frac{\partial u}{\partial y} = 4y + x - 2, \quad
        \frac{\partial u}{\partial z} = 6z - 6.
    $$

    在点 $(1,1,-1)$ 处计算得：
    $$
        \left.\frac{\partial u}{\partial x}\right|_{(1,1,-1)} = 6, \quad
        \left.\frac{\partial u}{\partial y}\right|_{(1,1,-1)} = 3, \quad
        \left.\frac{\partial u}{\partial z}\right|_{(1,1,-1)} = -12.
    $$

    因此梯度为
    $$
        \nabla u|_{(1,1,-1)} = (6, 3, -12).
    $$

    最大方向微商为
    $$
        \sqrt{6^2 + 3^2 + 12^2} = \sqrt{189}.
    $$
\end{sol}


26. 设 $z = f(xy)$, $f$ 为可微函数. 证明 $x\frac{\partial z}{\partial x} - y\frac{\partial z}{\partial y} = 0$.
\begin{pf}
    计算偏导数得
    $$
        \frac{\partial z}{\partial x} = y f'(xy), \quad
        \frac{\partial z}{\partial y} = x f'(xy).
    $$
    因此
    $$
        x\frac{\partial z}{\partial x} - y\frac{\partial z}{\partial y}
        = xy f'(xy) - xy f'(xy)
        = 0.
    $$

    Q.E.D.
\end{pf}


28. 证明函数 $u = \varphi(x - at) + \psi(x + at)$ 满足波动方程
$$
    \frac{\partial^2 u}{\partial t^2} = a^2 \frac{\partial^2 u}{\partial x^2}.
$$
其中 $\varphi, \psi$ 有连续的二阶微商.
\begin{pf}
    首先计算关于 $t$ 的偏导数:
    $$
        \frac{\partial u}{\partial t} = -a \varphi'(x - at) + a \psi'(x + at),
    $$
    $$
        \frac{\partial^2 u}{\partial t^2}
        = \frac{\partial}{\partial t} \left( \frac{\partial u}{\partial t} \right)
        = a^2 \varphi''(x - at) + a^2 \psi''(x + at).
    $$

    再计算关于 $x$ 的偏导数:
    $$
        \frac{\partial u}{\partial x} = \varphi'(x - at) + \psi'(x + at),
    $$
    $$
        \frac{\partial^2 u}{\partial x^2}
        = \frac{\partial}{\partial x} \left( \frac{\partial u}{\partial x} \right)
        = \varphi''(x - at) + \psi''(x + at).
    $$

    因此得到
    $$
        \frac{\partial^2 u}{\partial t^2} = a^2 \frac{\partial^2 u}{\partial x^2}.
    $$

    Q.E.D.
\end{pf}

\section*{第9章综合习题}

1. 设 $a_1,a_2,\dots,a_n$ 是非零常数. $f(x_1,x_2,\dots,x_n)$ 在 $\mathbb{R}^n$ 上可微. 求证: 存在 $\mathbb{R}$ 上一元可微函数 $F(s)$ 使得 $f(x_1,x_2,\dots,x_n) = F(a_1x_1 + a_2x_2 + \dots + a_nx_n)$ 的充分必要条件是
$$
    a_j\frac{\partial f}{\partial x_i} = a_i\frac{\partial f}{\partial x_j},\quad i,j = 1,2,\dots,n.
$$
\begin{pf}
    令$s=a_1x_1 + a_2x_2 + \dots + a_nx_n$
    \begin{itemize}
        \item 必要性

              因为存在 $\mathbb{R}$ 上一元可微函数 $F(s)$ 使得 $f(x_1,x_2,\dots,x_n) = F(a_1x_1 + a_2x_2 + \dots + a_nx_n)$，对二者求微分，得$\frac{\partial f}{\partial x_i}=a_iF'(s)$

              所以$a_j\frac{\partial f}{\partial x_i}=a_ja_iF'(s)=a_ia_jF'(s)=a_i\frac{\partial f}{\partial x_j}$

        \item 充分性

              因为$a_j\frac{\partial f}{\partial x_i} = a_i\frac{\partial f}{\partial x_j},\quad i,j = 1,2,\dots,n$，可以得到$\frac{1}{a_i}\frac{\partial f}{\partial x_i}$与$i$无关

              令$\varphi(x_1,\dots,x_n)=\frac{1}{a_i}\frac{\partial f}{\partial x_i}$

              所以$\frac{\partial f}{\partial x_i}=a_i\varphi$

              所以$df=\sum_{i=1}^na_i\varphi dx_i=\varphi ds$
    \end{itemize}
\end{pf}

\end{document}